\documentclass[a4paper]{article}

\usepackage[spanish]{babel}
\usepackage{graphicx}
\usepackage{amsmath, amssymb}
\usepackage[margin=2cm]{geometry}
\usepackage{fancyhdr}
\usepackage{enumerate}
\usepackage[shortlabels]{enumitem}
\usepackage{parskip}
\usepackage[most]{tcolorbox}
\usepackage[hidelinks]{hyperref}
\usepackage{float}


% cabecera
\pagestyle{fancy}
\fancyhead[l]{Daniel Soto}
\fancyhead[c]{Comando útiles Git-GitHub}
\fancyhead[r]{\today}
\fancyfoot[c]{\thepage}
\renewcommand{\headrulewidth}{0.2pt} % linea horizontal


\begin{document}

\textbf{Cómo instalar Git en Windows}

\begin{itemize}
	\item Presionar \texttt{Windows + R}
	\item Escribir \texttt{cmd}
	\item Presionar \texttt{Enter}
	\item Se va abrir la "Terminal" de Windows
	\item Escribir el siguiente comando y ejecutarlo (presionar \texttt{Enter}) \texttt{winget install Git.Git}
\end{itemize}


\textbf{Verificación de la instalación}

\begin{verbatim}
git --version
\end{verbatim}

\textbf{Configuración de Usuario y Email:} en el nombre puede poner lo que quiera, en el email, debe poner exactamente la dirección de correo electrónico de tu cuenta de GitHub

\begin{verbatim}
git config --global user.name "Tu Nombre"
git config --global user.email "tuemail@ejemplo.com"
\end{verbatim}

\textbf{Verificación de la configuración}

\begin{verbatim}
git config --list
\end{verbatim}

Ahora deberías ver como Output tu información.

\end{document}
